% Resume des commandes du robot (bibliotheque LibRobot de OCaml : open LibRobot)
%
% \apxCmd{...} ecrit la commande dans la premiere colonne et retient sa
% largeur ; \apxDef, au debut de la seconde colonne, decale le texte quand
% la commande deborde de la premiere colonne.
\newsavebox{\apxBoite}
\newcommand{\apxCmd}[1]{\global\setbox\apxBoite\hbox{\idt #1}\usebox{\apxBoite}}
\newcommand{\apxDef}{\apxFontSize\ifdim\wd\apxBoite>1.2cm
  \hspace*{\dimexpr\wd\apxBoite-1.2cm\relax}\fi}
% touches du clavier
\newcommand{\apxTouche}[1]{\mbox{\begin{picture}(12,12)(0, 2)
\put(0,0){\framebox(12,12)[lb]{}}
\put(1.5,2){\makebox(0,0)[lb]{\tiny #1}}
\end{picture}\hspace{0.3ex}}}
\newcommand{\apxEsc}{\mbox{\begin{picture}(14,12)(0, 2)
\put(0,0){\framebox(14,12)[lb]{}}
\put(1.5,2){\makebox(0,0)[lb]{\tiny Esc}}
\end{picture}\hspace{0.3ex}}}
\newcommand{\apxArret}{\apxEsc, \apxTouche{\S} ou \apxTouche{\#}}
\newcommand{\apxEntree}{\framebox[11.mm][l]{$\hookleftarrow$~~~}}
\newcommand{\apxEffacement}{\framebox[11.mm][l]{$\leftarrow$~~~}}
% chaines et caracteres de OCaml
\newcommand{\apxChaine}[1]{{\cha \char34{}#1\char34{}}}
\newcommand{\apxCar}[1]{{\cha \textquotesingle{}#1\textquotesingle{}}}
\newcommand{\apxRobot}{{\normalsize \rm\robot}}

\noindent{\apxFontSize Les commandes du \apxRobot\ sont
celles de la bibliothèque {\idt LibRobot} (un programme commence par
{\mcl open} {\idt LibRobot}). Une commande sans paramètre s'écrit suivie de
{\idt ()}, comme {\idt av ()}~; les paramètres sont séparés par des espaces
et ceux qui sont des expressions sont mis entre parenthèses, comme
{\idt ec 2} ou {\idt tl (-3.)}. Les commandes qui prennent deux coordonnées
les donnent entre parenthèses, séparées par une virgule, comme
{\idt deplacer (2, 3)}. Un nombre réel (de type {\idt float}) s'écrit avec
un point, comme {\idt lg 20.}, un entier (de type {\idt int}) sans point.
Les fonctions {\idt val}\ldots\ rendent la valeur courante d'un réglage du
\apxRobot.\par}
\vspace{3mm}

\noindent\begin{tabular}{@{}|p{1cm}p{11.9cm}|@{}} \hline
\multicolumn{2}{|l|}{Grille des points d'arrêt du {\normalsize \robot}~:}  \\ \hline \hline
\apxCmd{dg ()} & \apxDef \underline{d}essine la \underline{g}rille des points d'arrêt du \apxRobot\ (le
dessin de la grille peut être interrompu par \apxArret). \\
\apxCmd{lg $n$} & \apxDef définit la nouvelle \underline{l}ongueur $n$ (un réel
$n\:\geq\:0.$) des côtés d'un carreau de la \underline{g}rille ($12. \times 12.$ pixels
par défaut). \\
\apxCmd{lgX $n$} & \apxDef définit la nouvelle \underline{l}ongueur $n$ (réelle) du côté
horizontal (\underline{X}) d'un carreau de la \underline{g}rille ($12.$ pixels par défaut). \\
\apxCmd{lgY $n$} & \apxDef définit la nouvelle \underline{l}ongueur $n$ (réelle) du côté
vertical (\underline{Y}) d'un carreau de la \underline{g}rille ($12.$ pixels par défaut). \\
\apxCmd{lgRT ($r$, $t$)} & \apxDef définit les nouvelles \underline{l}ongueurs des
côtés d'un carreau de la \underline{g}rille par les coordonnées polaires \underline{$r$} et
\underline{$t$} (deux réels~: la longueur $r\:\geq\:0.$ de la diagonale du carreau est
exprimée en pixels et l'angle $t$ que fait cette diagonale avec le côté horizontal est
exprimé en degrés, en général $0.\:\leq\:t\:\leq\:90.$).
\\ \hline
\apxCmd{valLgX ()} & \apxDef est égal à la \underline{val}eur (de type {\idt float}) de la
\underline{l}on\underline{g}ueur du côté horizontal (\underline{X}) d'un carreau de la grille
(comptée en pixels). \\
\apxCmd{valLgY ()} & \apxDef est égal à la \underline{val}eur (de type {\idt float}) de la
\underline{l}on\underline{g}ueur du côté vertical (\underline{Y}) d'un carreau de la grille
(comptée en pixels). \\
\apxCmd{valLgR ()} & \apxDef est égal à la \underline{val}eur (de type {\idt float}) de la
\underline{l}on\underline{g}ueur de la diagonale (\underline{R}) d'un carreau de la grille
(comptée en pixels). \\
\apxCmd{valLgT ()} & \apxDef est égal à la \underline{val}eur (de type {\idt float}) de
\underline{l}'an\underline{g}le (\underline{T}) de la diagonale d'un carreau de la grille avec
son côté horizontal (compté en degrés). \\
\hline \end{tabular}

\noindent\begin{tabular}{@{}|p{1cm}p{11.9cm}|@{}} \hline
\apxCmd{valX ()} & \apxDef est égal à la \underline{val}eur (de type {\idt float}) de
l'abscisse absolue (\underline{X}) du \apxRobot\ (comptée en pixels à partir du bord gauche
de la fenêtre). \\
\apxCmd{valY ()} & \apxDef est égal à la \underline{val}eur (de type {\idt float}) de
l'ordonnée absolue (\underline{Y}) du \apxRobot\ (comptée en pixels à partir du bas de la
fenêtre), de sorte qu'après {\idt {\mcl let} x = round (valX ()) {\mcl and} y = round (valY
())}, plus tard {\idt lc ()$;$ tracer (x, y)} replace le \apxRobot\ à cette place sur
l'écran. \\ \hline
\end{tabular}

\noindent\begin{tabular}{@{}|p{1cm}p{11.9cm}|@{}} \hline
\multicolumn{2}{|l|}{Pivotements relatifs sur place~:}  \\ \hline \hline
\apxCmd{phtd ()} & \apxDef fait \underline{p}ivoter le \apxRobot\ sur place d'un
\underline{h}uitième de \underline{t}our à \underline{d}roite. \\
\apxCmd{pqtd ()} & \apxDef fait \underline{p}ivoter le \apxRobot\ sur place d'un
\underline{q}uart de \underline{t}our à \underline{d}roite. \\
\apxCmd{p3htd ()} & \apxDef fait \underline{p}ivoter le \apxRobot\ sur place de
\underline{3} \underline{h}uitièmes de \underline{t}our à \underline{d}roite. \\
\apxCmd{phtg ()} & \apxDef fait \underline{p}ivoter le \apxRobot\ sur place d'un
\underline{h}uitième de \underline{t}our à \underline{g}auche. \\
\apxCmd{pqtg ()} & \apxDef fait \underline{p}ivoter le \apxRobot\ sur place d'un
\underline{q}uart de \underline{t}our à \underline{g}auche. \\
\apxCmd{p3htg ()} & \apxDef fait \underline{p}ivoter le \apxRobot\ sur place de
\underline{3} \underline{h}uitièmes de \underline{t}our à \underline{g}auche. \\
\apxCmd{pdt ()} & \apxDef fait \underline{p}ivoter le \apxRobot\ sur place d'un
\underline{d}emi-\underline{t}our. \\ \hline
\end{tabular}

\noindent\begin{tabular}{@{}|p{1cm}p{11.9cm}|@{}} \hline
\multicolumn{2}{|l|}{Pivotements absolus sur place~: }  \\ \hline \hline
\apxCmd{pn ()} & \apxDef fait \underline{p}ivoter le \apxRobot\ sur place vers le nord (\underline{N}). \\
\apxCmd{pne ()} & \apxDef fait \underline{p}ivoter le \apxRobot\ sur place vers le nord-est (\underline{NE}). \\
\apxCmd{pe ()} & \apxDef fait \underline{p}ivoter le \apxRobot\ sur place vers l'est (\underline{E}). \\
\apxCmd{pse ()} & \apxDef fait \underline{p}ivoter le \apxRobot\ sur place vers le sud-est (\underline{SE}). \\
\apxCmd{ps ()} & \apxDef fait \underline{p}ivoter le \apxRobot\ sur place vers le sud (\underline{S}). \\
\apxCmd{psw ()} & \apxDef fait \underline{p}ivoter le \apxRobot\ sur place vers le sud-ouest (\underline{SW}). \\
\apxCmd{pw ()} & \apxDef fait \underline{p}ivoter le \apxRobot\ sur place vers l'ouest (\underline{W}). \\
\apxCmd{pnw ()} & \apxDef fait \underline{p}ivoter le \apxRobot\ sur place vers le nord-ouest (\underline{NW}). \\
\apxCmd{pivoter $o$} & \apxDef fait \underline{pivoter} le \apxRobot\ dans
l'\underline{o}rientation indiquée par la valeur de l'entier $o$ (entre 0 et 7) égale à
{\idt nord} $=$ 0, {\idt nordEst} $=$ 1, {\idt est} $=$ 2, {\idt sudEst} $=$ 3,
{\idt sud} $=$ 4, {\idt sudOuest} $=$ 5, {\idt ouest} $=$ 6 ou {\idt nordOuest} $=$ 7. \\
\apxCmd{rt $t$} & \apxDef change le repère et l'orientation du \apxRobot\ pour qu'il
effectue une \underline{r}o\underline{t}ation sur place d'un angle de $t$ degrés ($t$ réel,
dans le sens trigonométrique si $t > 0.$ et dans le sens des aiguilles d'une montre si
$t < 0.$, comme {\idt rt 15.} ou {\idt rt (-30.)}). \\ \hline
\apxCmd{valOrientation ()} & \apxDef est égal à la \underline{val}eur (de type {\idt int})
de l'\underline{orientation} du \apxRobot\ codée par {\idt nord} $=$ 0, {\idt nordEst} $=$ 1,
{\idt est} $=$ 2, {\idt sudEst} $=$ 3, {\idt sud} $=$ 4, {\idt sudOuest} $=$ 5,
{\idt ouest} $=$ 6 et {\idt nordOuest} $=$ 7. Les directions {\idt nord}, {\idt est},
{\idt sud} et {\idt ouest} sont parallèles aux bords de l'écran, le {\idt nord} en haut.
Les directions {\idt nordEst}, {\idt sudEst}, {\idt sudOuest} et {\idt nordOuest}
correspondent à l'orientation des diagonales d'un carreau de la grille.\\ \hline
\end{tabular}

\noindent\begin{tabular}{@{}|p{1cm}p{11.9cm}|@{}} \hline
\multicolumn{2}{|l|}{Crayon~: }  \\ \hline \hline
\apxCmd{bc ()} & \apxDef \underline{b}aisse le \underline{c}rayon pour dessiner lors des déplacements. \\
\apxCmd{lc ()} & \apxDef \underline{l}ève le \underline{c}rayon pour éviter de dessiner lors des déplacements. \\
\apxCmd{dc ()} & \apxDef \underline{d}essine une petite \underline{c}roix (de l'épaisseur du
crayon), sans déplacer le \apxRobot. \\
\apxCmd{dp ()} & \apxDef \underline{d}essine un \underline{p}oint (de l'épaisseur du
crayon), sans déplacer le \apxRobot. \\
\apxCmd{ec $n$} & \apxDef définit la nouvelle \underline{é}paisseur $n$ (un entier) du
\underline{c}rayon (1 pixel par défaut). \\ \hline
\apxCmd{valEc ()} & \apxDef est égal à la \underline{val}eur (de type {\idt int}) de
l'\underline{é}paisseur du \underline{c}rayon. \\
\apxCmd{valBc ()} & \apxDef est égal à la \underline{val}eur booléenne {\idt true} (vrai) si
et seulement si le \underline{c}rayon est \underline{b}aissé, sinon à {\idt false} (faux). \\ \hline
\end{tabular}

\noindent\begin{tabular}{@{}|p{1cm}p{11.9cm}|@{}} \hline
\multicolumn{2}{|l|}{Déplacements (avec ou sans tracé selon que le crayon est levé ou baissé)~: }  \\ \hline \hline
\apxCmd{av ()} & \apxDef fait \underline{av}ancer le \apxRobot\ tout droit jusqu'au
prochain point d'arrêt sur la grille. \\
\apxCmd{vd ()} & \apxDef fait tourner le \apxRobot\ par un \underline{v}irage d'un quart
de cercle à \underline{d}roite. \\
\apxCmd{vg ()} & \apxDef fait tourner le \apxRobot\ par un \underline{v}irage d'un quart
de cercle à \underline{g}auche. \\
\apxCmd{avf $n$} & \apxDef fait \underline{av}ancer le \apxRobot\ tout droit $n$
\underline{f}ois ($n$ entier~; sans rien faire si $n = 0$ et en reculant si $n < 0$, comme
{\idt avf (-2)}). \\
\apxCmd{tl $l$} & \apxDef change le repère (mais pas l'angle des directions NE-E) et fait
avancer le \apxRobot\ par une \underline{t}rans\underline{l}ation de $l$ pixels ($l$ réel)
dans la direction où il se trouve (si $l > 0.$) ou dans la direction opposée (si $l < 0.$).
\\
\apxCmd{tracerReel ($x$, $y$)} & \apxDef \underline{trace} un segment de droite (si le
crayon est baissé) de la position courante du \apxRobot\ au point de coordonnées réelles
$(x, y)$ du repère cartésien de la grille (voir {\idt deplacer}) et déplace le \apxRobot\ en
ce point. \\
\apxCmd{tracer ($x$, $y$)} & \apxDef fait de même jusqu'au point de coordonnées entières
$(x, y)$ comptées en pixels à partir du coin inférieur gauche de la fenêtre (comme
{\idt valX ()} et {\idt valY ()}).\\ \hline
\apxCmd{valLgTl ()} & \apxDef est égal à la \underline{val}eur (de type {\idt float}) de la
\underline{l}on\underline{g}ueur du segment tracé par la dernière
\underline{t}rans\underline{l}ation par {\idt av ()} ou {\idt tl} (comptée en pixels). \\ \hline
\end{tabular}

\noindent\begin{tabular}{@{}|p{1cm}p{11.9cm}|@{}} \hline
\multicolumn{2}{|l|}{Positionnement au centre et sur les bords du cadre~: }  \\ \hline \hline
\apxCmd{ce ()} & \apxDef place le \apxRobot\ au \underline{ce}ntre de {\em l'écran}, qui
est aussi {\em le centre de la grille}, sans changer l'orientation du \apxRobot\ et sans
rien dessiner. \\
\apxCmd{eh ()} & \apxDef place le \apxRobot\ \underline{e}n \underline{h}aut du cadre,
sans changer sa position horizontale, ni son orientation et sans rien dessiner. \\
\apxCmd{eb ()} & \apxDef place le \apxRobot\ \underline{e}n \underline{b}as du cadre, sans
changer sa position horizontale, ni son orientation et sans rien dessiner. \\
\apxCmd{ag ()} & \apxDef place le \apxRobot\ \underline{à} \underline{g}auche du cadre,
sans changer sa position verticale, ni son orientation et sans rien dessiner. \\
\apxCmd{ad ()} & \apxDef place le \apxRobot\ \underline{à} \underline{d}roite du cadre,
sans changer sa position verticale, ni son orientation et sans rien dessiner. \\ \hline
\end{tabular}

\noindent\begin{tabular}{@{}|p{1cm}p{11.9cm}|@{}} \hline
\multicolumn{2}{|l|}{Repère cartésien~: }  \\ \hline \hline
\apxCmd{translater ($x$, $y$)} & \apxDef \underline{translate} le \apxRobot\ de $x$
carreaux de la grille horizontalement (à droite si $x > 0$ et à gauche si $x < 0$) et $y$
carreaux de la grille verticalement (en haut si $y > 0$ et en bas si $y < 0$), sans changer
son orientation et sans rien dessiner ($x$ et $y$ entiers). \\
\apxCmd{deplacer ($x$, $y$)} & \apxDef \underline{déplace} le \apxRobot\ au point de
coordonnées entières $(x, y)$ sur la grille, sans changer son orientation ni rien dessiner,
l'origine du repère cartésien étant placée au centre de la grille. \\
\apxCmd{deplacerReel ($x$, $y$)} & \apxDef fait de même avec des coordonnées
\underline{réel}les $x$ et $y$ (de type {\idt float}), comme {\idt deplacerReel (1.5, -2.)}. \\ \hline
\end{tabular}

\noindent\begin{tabular}{@{}|p{1cm}p{11.9cm}|@{}} \hline
\multicolumn{2}{|l|}{Test booléen de l'orientation du {\normalsize \rm\robot}~:  }  \\ \hline \hline
\apxCmd{tn ()} & \apxDef \underline{t}este si le \apxRobot\ est orienté au nord (\underline{N}). \\
\apxCmd{tne ()} & \apxDef \underline{t}este si le \apxRobot\ est orienté au nord-est (\underline{NE}). \\
\apxCmd{te ()} & \apxDef \underline{t}este si le \apxRobot\ est orienté à l'est (\underline{E}).\\
\apxCmd{tse ()} & \apxDef \underline{t}este si le \apxRobot\ est orienté au sud-est (\underline{SE}).\\
\apxCmd{ts ()} & \apxDef \underline{t}este si le \apxRobot\ est orienté au sud (\underline{S}).\\
\apxCmd{tsw ()} & \apxDef \underline{t}este si le \apxRobot\ est orienté au sud-ouest (\underline{SW}).\\
\apxCmd{tw ()} & \apxDef \underline{t}este si le \apxRobot\ est orienté à l'ouest (\underline{W}).\\
\apxCmd{tnw ()} & \apxDef \underline{t}este si le \apxRobot\ est orienté au nord-ouest (\underline{NW}). \\ \hline
\end{tabular}

\noindent\begin{tabular}{@{}|p{1cm}p{11.9cm}|@{}} \hline
\multicolumn{2}{|l|}{Test booléen de la position du {\normalsize \rm\robot}~:  }  \\ \hline \hline
\apxCmd{tb ()} & \apxDef \underline{t}este si le \apxRobot\ est au \underline{b}ord du cadre
(un déplacement du \apxRobot\ en avant par {\idt av ()} le faisant sortir du cadre). \\ \hline
\end{tabular}

\noindent\begin{tabular}{@{}|p{1cm}p{11.9cm}|@{}} \hline
\multicolumn{2}{|l|}{Temps~:  }  \\ \hline \hline
\apxCmd{heure h mi s} & \apxDef met dans les variables entières {\idt h}, {\idt mi} et
{\idt s} (déclarées par {\idt {\mcl let} h = ref 0}, \ldots) l'heure (entre 0 et 23), la
minute (entre 0 et 59) et la seconde (entre 0 et 59) présentes, qui sont alors les valeurs
de {\idt !h}, {\idt !mi} et {\idt !s}. \\
\apxCmd{date j mo a} & \apxDef met de même dans les variables entières {\idt j}, {\idt mo}
et {\idt a} le jour (entre 1 et 31), le mois (entre 1 et 12) et l'année courante. \\ \hline
\apxCmd{delai $t$} & \apxDef attend pendant un \underline{délai} de $t$ secondes ($t$ de type
{\idt float}, comme {\idt delai 0.5}). \\
\apxCmd{tickCount ()} & \apxDef est égal au nombre entier de soixantièmes de seconde
écoulés (pour mesurer une durée par différence de deux valeurs). \\ \hline
\end{tabular}

\noindent\begin{tabular}{@{}|p{1cm}p{11.9cm}|@{}} \hline
\multicolumn{2}{|l|}{Interruption de l'exécution du programme~:  }  \\ \hline \hline
\apxCmd{\rm\apxArret} & \apxDef tapée pendant que le \apxRobot\ se déplace (par
{\idt av ()}, {\idt vd ()}, {\idt vg ()}, \ldots) ou s'exécute pas à pas, l'une de ces
touches arrête le programme~; la fenêtre reste affichée jusqu'à ce que l'on tape une touche
ou que l'on clique. \\ \hline
\end{tabular}

\noindent\begin{tabular}{@{}|p{1cm}p{11.9cm}|@{}} \hline
\multicolumn{2}{|l|}{ Mise au point du programme~:  }  \\ \hline \hline
\apxCmd{st ()} & \apxDef \underline{st}oppe le \apxRobot\ et rappelle dans la ligne du bas
de la fenêtre les touches utilisables~: \apxTouche{$+$} augmente et \apxTouche{$-$}
diminue la vitesse du \apxRobot\ (entre 0 et 10, la vitesse 0 faisant passer en mode pas à
pas), \apxEsc\ ou \apxTouche{\#} ferme la fenêtre et termine le programme, toute autre
touche fait redessiner tout le dessin depuis le début, à la vitesse choisie. Un programme
qui se termine sans {\idt st ()} attend que l'on tape une touche ou que l'on clique avant de
fermer la fenêtre. \\
\apxCmd{pp ()} & \apxDef exécute la suite du programme en mode \underline{p}as à
\underline{p}as, à la vitesse 0~: le \apxRobot\ attend que l'on tape une touche avant
chaque commande (\apxArret\ arrêtant le programme). \\
\apxCmd{ex ()} & \apxDef exécute la suite du programme en mode \underline{ex}écution
normale, à la vitesse 10. \\
\apxCmd{vt $v$} & \apxDef exécute la suite du programme à la \underline{v}i\underline{t}esse
$v$ (un entier~; en pas à pas si $v \leq 0$, avec un maximum de 10 si $v > 10$). \\
\apxCmd{montrerRobot ()} & \apxDef rend le \apxRobot\ visible sur l'écran pendant le dessin
(option par défaut). \\
\apxCmd{cacherRobot ()} & \apxDef rend le \apxRobot\ invisible sur l'écran pendant le dessin. \\
\apxCmd{interdireSorties ()} & \apxDef \underline{interdit} les \underline{sorties} du
\apxRobot\ en dehors du cadre dans la suite du programme (option par défaut~: une commande
qui ferait sortir le \apxRobot\ du cadre arrête alors le programme sur une erreur). \\
\apxCmd{autoriserSorties ()} & \apxDef \underline{autorise} les \underline{sorties} du
\apxRobot\ en dehors du cadre dans la suite de l'exécution du programme. \\ \hline
\apxCmd{valVt ()} & \apxDef est égal à la \underline{val}eur (de type {\idt int}) de la
\underline{v}i\underline{t}esse du \apxRobot\ (entre 0 (mode pas à pas) et 10). \\
\apxCmd{valVisible ()} & \apxDef est égal à la valeur booléenne {\idt true} (respectivement
{\idt false}) si le \apxRobot\ est \underline{visible} (respectivement invisible) sur l'écran. \\
\apxCmd{valAutoriserSorties ()} & \apxDef est égal à la valeur booléenne {\idt true}
(respectivement {\idt false}) si les \underline{sorties} en dehors du cadre sont
\underline{autorisées} (respectivement interdites). \\ \hline
\end{tabular}

\noindent\begin{tabular}{@{}|p{1cm}p{11.9cm}|@{}} \hline
%
Couleur~: & \\ \hline \hline
%
\multicolumn{2}{|p{12.9cm}|}{Les couleurs sont des entiers~: {\idt blanc} $= 0$,
{\idt jaune} $= 1$, {\idt bleu} $= 2$, {\idt violet} $= 3$, {\idt vert} $= 4$,
{\idt rouge} $= 5$, {\idt indigo} $= 6$ et {\idt noir} $= 7$.}\\\hline \hline
%
\apxCmd{cc $c$} & \apxDef définit la \underline{c}ouleur $c$ du \underline{c}rayon,
comme {\idt cc rouge}. \\
\apxCmd{peindre ()} & \apxDef \underline{peint} de la couleur du crayon la région
d'une même couleur (blanche en général), limitée par les traits du dessin, où se trouve le
\apxRobot. \\
\apxCmd{cf $c$} & \apxDef efface complètement le dessin puis remplit l'écran de la
\underline{c}ouleur de \underline{f}ond $c$. Le \apxRobot\ est placé au centre de l'écran
dans les conditions initiales (visible, tourné vers le nord, sorties hors du cadre
interdites, crayon noir d'épaisseur 1 baissé et exécution normale à la vitesse 10)~; les
dimensions de la grille ne changent pas. \\ \hline
\apxCmd{valCf ()} & \apxDef est égal à la \underline{val}eur (de type {\idt int}) de la
\underline{c}ouleur du \underline{f}ond. \\
\apxCmd{valCc ()} & \apxDef est égal à la \underline{val}eur (de type {\idt int}) de la
\underline{c}ouleur du \underline{c}rayon. \\
\apxCmd{ecranNoirEtBlanc ()} & \apxDef est égal à {\idt true} si et seulement si l'écran est
noir et blanc, sinon à {\idt false} (c'est toujours {\idt false} sur les écrans actuels). \\ \hline
\end{tabular}

\noindent\begin{tabular}{@{}|p{1cm}p{11.9cm}|@{}} \hline
\multicolumn{2}{|l|}{Écriture dans la fenêtre de dessin du {\normalsize \rm\robot}~:  }  \\ \hline \hline
\apxCmd{ecrire $c$} & \apxDef écrit la chaîne de caractères $c$ (entre guillemets
\apxChaine{} et \apxChaine{}, où tout guillemet doit être précédé d'une barre oblique
inverse~: {\cha \textbackslash\char34}) à droite du \apxRobot\ (ou à sa gauche s'il est
tourné vers l'ouest), sans changer sa position ni celle de son crayon. \\
\apxCmd{entierEnChaine $i$} & \apxDef est la chaîne de caractères représentant l'entier
$i$ en base 10 (comme {\idt string\_of\_int $i$}). \\
\apxCmd{reelEnChaine $r$} & \apxDef est la chaîne de caractères représentant le réel $r$
sous la forme {\idt $d.dddd$e$\pm ee$} où $d$ est un chiffre de la mantisse, $e$ est un
chiffre de l'exposant et {\idt e} se lit «~fois 10 à la puissance~». \\
\apxCmd{booleenEnChaine $b$} & \apxDef est la chaîne de caractères
(\apxChaine{Vrai} ou \apxChaine{Faux}) re\-pré\-sen\-tant le booléen $b$.\\
\apxCmd{concat [$c_1$$;$ $c_2$$;$ \ldots$;$ $c_n$]} & \apxDef est la chaîne de
caractères obtenue en \underline{concat}énant les chaînes $c_1$, $c_2$, \ldots, $c_n$ les
unes derrière les autres (comme $c_1$ {\idt \char94} $c_2$ {\idt \char94} \ldots\
{\idt \char94} $c_n$). \\ \hline
\apxCmd{policeTexte $p$} & \apxDef choisit la \underline{police} $p$ du
\underline{texte} ({\idt courier}, {\idt times}, {\idt helvetica}, \ldots). \\
\apxCmd{tailleTexte $t$} & \apxDef choisit la \underline{taille} $t$ (un entier) du
\underline{texte}. \\
\apxCmd{styleTexte $s$} & \apxDef choisit le \underline{style} du \underline{texte}
({\idt normal}, {\idt gras}, {\idt italique}, {\idt souligne}, \ldots). \\
\multicolumn{2}{|p{12.9cm}|}{\apxFontSize Ces trois commandes sont sans effet visible~: la
police de la fenêtre du \apxRobot\ ne peut pas être changée.} \\ \hline
\apxCmd{valPoliceTexte ()} & \apxDef \underline{val}eur (de type {\idt string}) de la
\underline{police} courante du \underline{texte}. \\
\apxCmd{valTailleTexte ()} & \apxDef \underline{val}eur (de type {\idt int}) de la
\underline{taille} courante du \underline{texte}. \\ \hline
\end{tabular}

\noindent\begin{tabular}{@{}|p{1cm}p{11.9cm}|@{}} \hline
\multicolumn{2}{|l|}{Messages affichés dans la ligne du bas de la fenêtre du {\normalsize \rm\robot}~:  }  \\ \hline \hline
\apxCmd{message $c$} & \apxDef efface le \underline{message} de la ligne du bas de la
fenêtre du \apxRobot\ puis y écrit la chaîne de caractères $c$. \\
\apxCmd{marquerUnePause ()} & \apxDef \underline{marque} \underline{une}
\underline{pause} (en faisant clignoter le symbole $\hookleftarrow$ dans la ligne des
messages) en attendant que le bouton de la souris ou qu'une touche du clavier soit enfoncé.\\
\apxCmd{effacerMessage ()} & \apxDef \underline{efface} le \underline{message} de la
ligne du bas de la fenêtre du \apxRobot. \\ \hline
\apxCmd{lireCar ()} & \apxDef attend (en faisant clignoter le symbole $\hookleftarrow$
dans la ligne des messages) qu'une touche du clavier ou que le bouton de la souris soit
enfoncé, puis \underline{lit} et rend le \underline{car}actère (de type {\idt char})
correspondant à la touche enfoncée (ou le caractère \apxCar{0} si le bouton de la souris a
été enfoncé). \\ \hline
\apxCmd{mementoCommandesClavier} & \apxDef est une chaîne de caractères
rappelant les touches de pilotage interactif du \apxRobot\ (voir {\idt interaction}), à
afficher par {\idt message mementoCommandesClavier}. \\
\apxCmd{valMessage ()} & \apxDef est égal à la \underline{val}eur (de type {\idt string})
du \underline{message} inscrit dans la ligne du bas de la fenêtre du \apxRobot\ ou à la
chaîne vide \apxChaine{} s'il n'y a aucun message. \\ \hline
\end{tabular}
\vfil
\noindent\begin{tabular}{@{}|p{1cm}p{11.9cm}|@{}} \hline
\multicolumn{2}{|l|}{Lecture et écriture de texte~:  }  \\ \hline \hline
\multicolumn{2}{|p{12.9cm}|}{\apxFontSize Le texte écrit par {\idt write} et
{\idt writeln} s'affiche dans le terminal. Quand le programme attend une réponse
({\idt lireEntier ()}, \ldots), la question (la ligne en cours, précédée des deux
dernières lignes écrites) s'affiche dans la ligne du bas de la fenêtre du \apxRobot. La
réponse se tape dans la fenêtre du \apxRobot, où elle s'affiche à la suite de la question
(\apxEffacement\ pour corriger, \apxEntree\ pour terminer), ou bien dans le terminal.} \\ \hline
\apxCmd{write [$c_1$$;$ \ldots$;$ $c_n$]} & \apxDef écrit les chaînes de caractères
$c_1$, \ldots, $c_n$ les unes à la suite des autres, comme
{\idt write [\apxChaine{N = }$;$ entierEnChaine n]}. \\
\apxCmd{writeln [$c_1$$;$ \ldots$;$ $c_n$]} & \apxDef fait de même puis passe à la ligne
({\idt writeln []} passe seulement à la ligne). \\ \hline
\apxCmd{lireEntier ()} & \apxDef \underline{lit} la réponse de l'utilisateur, terminée par
\apxEntree, et rend l'\underline{entier} (de type {\idt int}) qu'elle représente. Si la
réponse n'est pas un nombre entier, la question \apxChaine{Entrer un nombre entier : } est
posée à nouveau. Sur un clavier français, les chiffres peuvent être tapés sans la touche
majuscule. \\
\apxCmd{lireReel ()} & \apxDef fait de même pour \underline{lire} un nombre
\underline{réel} (de type {\idt float}). \\
\apxCmd{lireChaine ()} & \apxDef \underline{lit} la réponse de l'utilisateur, terminée par
\apxEntree, et la rend comme une \underline{chaîne} de caractères (de type {\idt string}). \\
\apxCmd{lireCaractere ()} & \apxDef fait de même et rend le premier \underline{caractère}
de la réponse (de type {\idt char}). \\
\apxCmd{readln ()} & \apxDef attend que l'utilisateur tape \apxEntree\ (ce qu'il a tapé
avant est ignoré). \\
\apxCmd{effacerTexte ()} & \apxDef \underline{efface} le \underline{texte} du terminal. \\ \hline
\end{tabular}
\vfil
\noindent\begin{tabular}{@{}|p{1cm}p{11.9cm}|@{}} \hline
\multicolumn{2}{|l|}{Fonctions utiles~:  }  \\ \hline \hline
\apxCmd{pi} & \apxDef est le réel $\pi = 3.14159\ldots$ (de type {\idt float}). \\
\apxCmd{round $x$} & \apxDef est l'entier le plus proche du réel $x$, comme
{\idt round 2.5} $=$ 3. \\
\apxCmd{random ()} & \apxDef est un entier tiré au hasard entre $-32767$ et $32767$. \\ \hline
\end{tabular}
\vfil
\noindent\begin{tabular}{@{}|p{1cm}p{11.9cm}|@{}} \hline
\multicolumn{2}{|l|}{Sons~:  }  \\ \hline \hline
\apxCmd{bip ()} & \apxDef émet un son bref. \\
\apxCmd{biiip ()} & \apxDef émet un son long. \\ \hline
\end{tabular}
\vfil
\noindent\begin{tabular}{@{}|p{1cm}p{11.9cm}|@{}} \hline
\multicolumn{2}{|l|}{Musique~:  }  \\ \hline \hline
\apxCmd{metronome $n$} & \apxDef règle le \underline{métronome} à $n$ noires par minute
({\idt largo} $=$ 50, {\idt lento} $=$ 56, {\idt larghetto} $=$ 60, {\idt adagio} $=$ 70,
{\idt andante} $=$ 85, {\idt andantino} $=$ 95, {\idt moderato} $=$ 110,
{\idt allegretto} $=$ 115, {\idt allegro} $=$ 130, {\idt vivace} $=$ 150,
{\idt presto} $=$ 175, {\idt prestissimo} $=$ 200~; 100 par défaut). \\
\apxCmd{nuance $n$} & \apxDef règle la \underline{nuance} $n$ (entre 1 et 9~:
{\idt pianississimo} $=$ 1, {\idt pianissimo} $=$ 2, {\idt piano} $=$ 3,
{\idt mezza\_voce} $=$ 4, {\idt mezzo\_piano} $=$ 4, {\idt mezzo\_forte} $=$ 6,
{\idt forte} $=$ 7, {\idt fortissimo} $=$ 8, {\idt fortississimo} $=$ 9). \\
\apxCmd{note $h$ $d$} & \apxDef joue une \underline{note} de hauteur
$h = o * \mbox{\idt octave} + n + a$ (un entier) où $o = -3$, $-2$, \ldots, $3$ est l'octave
de la note (0 pour l'octave du do du milieu du piano, {\idt octave} $=$ 12), $n$ $=$
{\idt do0} ou {\idt d0} ou {\idt ut} $=$ 0, {\idt re} $=$ 2, {\idt mi} $=$ 4, {\idt fa} $=$ 5,
{\idt sol} $=$ 7, {\idt la} $=$ 9, {\idt si} $=$ 11 est le nom de la note, $a$ $=$
{\idt bemol} $= -1$, {\idt becarre} $= 0$, {\idt diese} $= 1$ est l'altération de la note
et $d$ (un réel) est la durée de la note, comptée en noires~: {\idt ronde} $= 4.$,
{\idt ronde\_pointee} $= 6.$, {\idt blanche} $= 2.$, {\idt blanche\_pointee} $= 3.$,
{\idt blanche\_en\_triolet} $= 4/3$, {\idt noire} $= 1.$, {\idt noire\_pointee} $= 1.5$,
{\idt noire\_en\_triolet} $= 2/3$, {\idt croche} $= 0.5$, {\idt croche\_pointee} $= 0.75$,
{\idt croche\_en\_triolet} $= 1/3$, {\idt double\_croche} $= 0.25$,
{\idt double\_croche\_pointee} $= 0.375$ ou {\idt triple\_croche} $= 0.125$, comme
{\idt note (octave + sol + diese) croche}. \\
\apxCmd{silence $d$} & \apxDef marque un \underline{silence} de durée $d$ $=$
{\idt pause} $= 4.$, {\idt demi\_pause\_pointee} $= 3.$, {\idt demi\_pause} $= 2.$,
{\idt soupir\_pointe} $= 1.5$, {\idt soupir} $= 1.$, {\idt demi\_soupir\_pointe} $= 0.75$,
{\idt demi\_soupir} $= 0.5$, {\idt quart\_de\_soupir} $= 0.25$ ou
{\idt huitieme\_de\_soupir} $= 0.125$. \\ \hline
\end{tabular}
\vfil
\noindent\begin{tabular}{@{}|p{1cm}p{11.9cm}|@{}} \hline
\multicolumn{2}{|l|}{Interaction avec l'utilisateur du {\normalsize \rm\robot}~:  }  \\ \hline \hline
\multicolumn{2}{|l|}{\idt interaction (coordonnees, pivotements, dessin, finSouris, sauver)}\\
& \apxFontSize efface l'écran puis laisse l'utilisateur déplacer le \apxRobot\ avec la souris,
en cliquant sur un point de la grille (si {\idt finSouris} est {\idt true}~: de carreau en carreau, bouton relâché), et éventuellement le piloter au clavier (les cinq
paramètres sont des booléens). L'interaction se termine par la touche \apxEntree\ (ou en
enfonçant le bouton de la souris si {\idt finSouris} est {\idt true})~; le programme
s'arrête alors comme avec {\idt st ()}.

\tiret Si {\idt coordonnees} est {\idt true}, les coordonnées du \apxRobot\ sont affichées
pendant l'interaction.

\tiret Si {\idt pivotements} est {\idt true}, l'utilisateur peut faire pivoter le
\apxRobot\ avec les touches \apxTouche{$-$} ({\idt pqtg}), \apxTouche{$+$} ({\idt pqtd}),
\apxTouche{g} ({\idt phtg}) et \apxTouche{d} ({\idt phtd}).

\tiret Si {\idt pivotements} et {\idt dessin} sont {\idt true}, il peut aussi dessiner avec
les touches \apxTouche{a} ({\idt av}), \apxTouche{$($} ({\idt vg}), \apxTouche{$)$}
({\idt vd}), \apxTouche{l} ({\idt lc}), \apxTouche{b} ({\idt bc}), \apxTouche{p}
({\idt dp}), \apxTouche{c} ({\idt dc}), \apxTouche{$.$} ({\idt dg}) et \apxTouche{$*$}
(changement de la couleur du crayon), en annulant la dernière commande avec
\apxEffacement.

\tiret Si {\idt dessin} et {\idt sauver} sont {\idt true}, le programme demande à la fin
s'il faut sauver le dessin~: en répondant \apxTouche{y}, on obtient un programme
{\idt mon\_dessin.ml} dont l'exécution reproduit le dessin.
\\ \hline
\multicolumn{2}{|l|}{\idt interactionTP x y c coordonnees pivotements dessin finSouris sauver}\\
& \apxFontSize laisse l'utilisateur déplacer le \apxRobot\ avec la souris et le piloter au
clavier comme {\idt interaction} (sauf \apxTouche{$*$} et \apxEffacement), mais sans
effacer l'écran, et continue ensuite le
programme. L'interaction se termine par la touche \apxEntree, par \apxEsc, par toute touche
qui n'est pas une commande du \apxRobot\ ou, si {\idt finSouris} est {\idt true}, en
enfonçant le bouton de la souris. À la fin, les variables entières {\idt x} et {\idt y}
(déclarées par {\idt {\mcl let} x = ref 0}, \ldots) contiennent le déplacement horizontal et
vertical du \apxRobot\ compté en carreaux de la grille, et la variable {\idt c} (déclarée
par {\idt {\mcl let} c = ref \apxCar{ }}) contient la dernière touche tapée (un espace si
l'interaction s'est terminée par un clic). Le paramètre {\idt sauver} est sans effet.
\\ \hline
\end{tabular}
\newpage                 % Si on termine le chapitre par une page blanche
\thispagestyle{empty}
